\subsection{\texorpdfstring{\(\mathcal{M}(\mathrm{X};\mathbf{C})\) 中的紧收敛}{M(X; C) 中的紧收敛}}

回顾一下，\(\mathcal{M}(\mathrm{X};\mathbf{C})\) 上的紧收敛拓扑就是在 \(\mathcal{K}(\mathrm{X};\mathbf{C})\) 的紧子集上的一致收敛拓扑。我们把 \(\mathcal{M}(\mathrm{X};\mathbf{C})\) 上在 \(\mathcal{K}(\mathrm{X};\mathbf{C})\) 的严格紧子集（No.~1）上的一致收敛拓扑称为严格紧收敛拓扑。

命题 17. --- 在空间 \(\mathcal{M}(\mathrm{X};\mathbf{C})\) 上，考虑下列诸拓扑：

\(\mathcal{T}_1\) ：关于一个全子集 \(\mathrm{T}\) （\(\mathcal{K}(\mathrm{X};\mathbf{C})\) 的）的逐点收敛拓扑；

\(\mathcal{T}_2\) ：淡拓扑；

\(\mathcal{T}_3\) ：严格紧收敛拓扑；

\(\mathcal{T}_4\) ：紧收敛拓扑。

这些拓扑中的每一个都比其后一个粗。此外：

\begin{enumerate}
\def\labelenumi{(\roman{enumi})}
\item
  对 \(\mathcal{T}_2\) 、 \(\mathcal{T}_3\) 与 \(\mathcal{T}_4\) 而言，有界集是相同的。
\item
  若 \(\mathbf{H}\) 是 \(\mathcal{M}(\mathrm{X};\mathbf{C})\) 的一个淡有界子集，则 \(\mathbf{H}\) 上由 \(\mathcal{T}_1, \mathcal{T}_2, \mathcal{T}_3, \mathcal{T}_4\) 诸拓扑诱导的拓扑是相同的。
\end{enumerate}

\(\mathcal{M}(X;C)\) 的淡有界子集 H 是等度连续的（No.~9, 命题 15），于是第一个断言由 TVS, III, §3, No.~7, 命题 9 得出，第二个断言由 GT, X, §2, No.~4, Th. 1 得出。

回顾一下，当 X 仿紧时，严格紧收敛拓扑与紧收敛拓扑一致（No.~1, 命题 2）。

命题 18. --- 在锥 \(\mathcal{M}_{+}(X)\) 上，由下列诸拓扑诱导的拓扑彼此重合：

\(\mathcal{T}_{1}\) ：关于 \(\mathcal{K}(X;C)\) 的一个在 \(\mathcal{K}(X;C)\) 中稠密且满足性质（P）的线性子空间 V 的逐点收敛拓扑（No.~7, 命题 9）；

\(\mathcal{T}_2\) ：淡拓扑；

\(\mathcal{T}_3\) ：严格紧收敛拓扑。

由于每个滤子都是比它细的诸超滤子之交（GT, I, §6, No.~4, 命题 7），只需证明：若 \(\mathfrak{U}\) 是 \(\mathcal{M}_{+}(\mathrm{X})\) 上的一个超滤子，测度 \(\mu_0\) 是它对拓扑 \(\mathcal{T}_1\) 的极限，则 \(\mu_0\) 也是它对 \(\mathcal{T}_3\) 的极限。设 \(\mathbf{K}\) 是 \(\mathbf{X}\) 的一个紧子集；按假设，存在一个函数 \(h \in \mathrm{V}\) ，它 \(\geqslant 0\) 于 \(\mathbf{X}\) 且取值 \(>0\) 于 \(\mathbf{K}\) ；由此可知每个函数 \(f \in \mathcal{K}(\mathrm{X},\mathrm{K};\mathbf{C})\) 都可写成 \(f = gh\) ，其中 \(g \in \mathcal{K}(\mathrm{X},\mathrm{K};\mathbf{C})\) ，并且若 \(c = \inf_{x \in \mathrm{K}} h(x) > 0\) ，则 \(\| g \| \leqslant c^{-1} \| f \|\) 。按假设，存在一个集 \(\mathrm{H}_0 \in \mathfrak{U}\) ，使得对每个测度 \(\mu \in \mathrm{H}_0\) ，

\[
0 \leqslant \mu (h) \leqslant \mu_ {0} (h) + 1 = b.
\]

因而，对每个函数 \(f \in \mathcal{K}(\mathrm{X};\mathbf{C})\) ，有

\[
| \langle f, h \cdot \mu \rangle | = | \langle h f, \mu \rangle | \leqslant \| f \| \cdot \mu (h) \leqslant b \| f \|
\]

对每个测度 \(\mu\in\mathrm{H}_{0}\) 成立；这就证明：测度 \(h\cdot\mu\) （\(\mu\) 跑遍 \(\mathrm{H}_{0}\) ）所成的集 H 是淡有界的。若 \(\mathfrak{U}_{0}\) 是 \(\mathfrak{U}\) 在 \(\mathrm{H}_{0}\) 上诱导的超滤子，则 \(\mathfrak{U}_{0}\) 在映射 \(\mu\mapsto h\cdot\mu\) 下的像是 H 上某个超滤子 \(\mathfrak{F}\) 的基；又因 H 对严格紧收敛拓扑是相对紧的（命题 17 与 No.~9, 命题 15），故 \(\mathfrak{F}\) 对该拓扑收敛于某个测度 \(\nu_{0}\) 。换言之，对任意 \(\varepsilon>0\) 以及 \(\mathcal{K}(\mathrm{X},\mathrm{K};\mathbf{C})\) 的任意紧子集 L ，存在 \(\mathrm{H}_{0}\) 的一个属于 \(\mathfrak{U}\) 的子集 N ，使得对每个函数 \(g\in\mathrm{L}\) 及属于 N 的任意一对测度 \(\mu,\mu'\) ，有 \(|\langle g,h\cdot\mu\rangle-\langle g,h\cdot\mu'\rangle|\leqslant\varepsilon\) ，即

\[
| \langle g h, \mu \rangle - \langle g h, \mu^ {\prime} \rangle | \leqslant \varepsilon .
\]

而上面我们已经看到，映射 \(g \mapsto gh\) 是 Banach 空间 \(\mathcal{K}(\mathrm{X},\mathrm{K};\mathbf{C})\) 的一个自同构。于是我们已证明 U 是 \(\mathcal{M}_{+}(\mathrm{X})\) 上对严格紧收敛拓扑的 Cauchy 滤子。从而它更是对淡收敛的 Cauchy 滤子，且 No.~9 的命题 14 表明它淡收敛于某个测度 \(\mu_{1}\) ；此外，由于 V 在 \(\mathcal{K}(\mathrm{X};\mathbf{C})\) 中稠密，假设蕴涵 \(\mu_{1} = \mu_{0}\) ；最后，由于 U 是对严格紧收敛拓扑的 Cauchy 滤子，它对该拓扑也收敛于 \(\mu_{0}\) （GT, X, §1, No.~5, 命题 5）。

证毕。

推论. --- 若 X 仿紧，则淡拓扑与紧收敛拓扑在 \(\mathcal{M}_{+}(X)\) 上诱导的拓扑重合。

然而，当 X 不是仿紧的时候，紧收敛拓扑与严格紧收敛拓扑在 \(\mathcal{M}_{+}(X)\) 上诱导的拓扑可能不同（习题 3）。

